% Source: source/普通高考/2013/2013山东理.pdf, page 1, question 13.
% pdfimages -list reports no embedded bitmap; the flowchart is vector/text artwork.
% Grid evidence: tmp/review_flowchart_2013_shandong_science/grid/page-1-grid100.png
% (100px coarse + 50px fine) and q13-block-grid50.png (50px coarse + 25px fine).
% Comparison crop: tmp/review_flowchart_2013_shandong_science/crop/q13_original_final.png,
% taken from the ungridded page render with a safety border.
% Nodes: 开始; 输入 ε (ε>0); F_0=1,F_1=2,n=1; F_1=F_0+F_1;
% F_0=F_1−F_0; n=n+1; 1/F_1≤ε; 输出 n; 结束.
% Directed edges: 开始→输入→初始化→F_1更新→F_0更新→n更新→判定;
% 判定(是)→输出→结束; 判定(否)→右回线→共享入口→F_1更新.
% The shared loop entry has one left-pointing return arrow and one downward
% arrow into F_1更新; all lines are solid and all node text is centered with
% local zero paragraph indentation.
\begin{tikzpicture}[
  >=Stealth,
  line width=0.55pt,
  line cap=round,
  line join=round,
  every node/.style={font=\normalsize,anchor=center,align=center,
    execute at begin node={\setlength{\parindent}{0pt}\noindent}},
  flowbox/.style={draw,minimum height=.68cm,inner xsep=4pt,inner ysep=2pt,
    anchor=center,align=center},
  terminal/.style={flowbox,rounded corners=.18cm,minimum width=1.35cm},
  io/.style={flowbox,shape=trapezium,
    trapezium left angle=72,trapezium right angle=108,
    minimum height=.78cm},
  flowarrow/.style={->,line width=0.55pt},
]
  \node[terminal,minimum height=.72cm] (start) at (0,9.00) {开始};
  \node[io,minimum width=2.60cm] (input) at (0,7.72)
    {输入 $\varepsilon$ ($\varepsilon>0$)};
  \node[flowbox,minimum width=4.05cm] (init) at (0,6.48)
    {$F_0=1,F_1=2,n=1$};
  \node[flowbox,minimum width=2.85cm] (sum1) at (0,5.18)
    {$F_1=F_0+F_1$};
  \node[flowbox,minimum width=2.85cm] (sum0) at (0,3.92)
    {$F_0=F_1-F_0$};
  \node[flowbox,minimum width=1.90cm] (inc) at (0,2.66)
    {$n=n+1$};
  \node[flowbox,shape=diamond,aspect=2.75,minimum width=4.10cm,
    minimum height=1.05cm] (test) at (0,1.28)
    {$\dfrac{1}{F_1}\leq\varepsilon$};
  \node[io,minimum width=1.90cm] (output) at (0,-0.25) {输出 $n$};
  \node[terminal,minimum height=.72cm] (finish) at (0,-1.55) {结束};

  \coordinate (loopjoin) at (0,5.80);
  \coordinate (loop-right) at (3.15,5.80);
  \coordinate (loop-bottom) at (3.15,1.28);

  \draw[flowarrow] (start.south) -- (input.north);
  \draw[flowarrow] (input.south) -- (init.north);
  \draw (init.south) -- (loopjoin);
  \draw[flowarrow] (loopjoin) -- (sum1.north);
  \draw[flowarrow] (sum1.south) -- (sum0.north);
  \draw[flowarrow] (sum0.south) -- (inc.north);
  \draw[flowarrow] (inc.south) -- (test.north);
  \draw[flowarrow] (test.south) -- node[right=2pt] {是} (output.north);
  \draw[flowarrow] (output.south) -- (finish.north);

  % 否 branch loops right and returns left into the shared entry.
  \draw (test.east) -- (loop-bottom) -- (loop-right);
  \draw[flowarrow] (loop-right) -- (loopjoin);
  \node[inner sep=0pt,anchor=west] at (2.35,2.00) {否};

  \path[use as bounding box] (-2.35,-2.00) rectangle (3.70,9.45);
\end{tikzpicture}
